\section{DALL-E：Text-to-Image 与 Zero-Shot Transformation}

DALL-E 把 text token 与 image token 放进统一 autoregressive model，学习 $p(\text{image}\mid\text{text})$。课堂用 interpolation、复杂 prompt 和 image-to-image examples 展示组合能力。示例很有说服力，也最容易被 cherry-pick；因此必须把 sample gallery 当定性证据，而不是完整分布统计。

\subsection{概念插值与组合}

本节先从概念插值观察文本条件如何移动视觉分布；读图时要同时检查连续性、对象身份和失败样本，而不是只看画面是否有趣。前一节的 iGPT 只给视觉前缀，这里则由语言指定目标属性，因此条件遵循和视觉真实性必须分开评价。

\lecturefigure{slide-25-dalle-interpolation.jpg}{DALL-E 通过文本条件在概念之间生成连续变化的图像样本。}{官方视频 29:10；Ramesh et al. 2021。}

\begin{knowledgebox}{读图：插值展示条件控制，不证明概念几何完美}
从猫到犬等 prompt variation 让输出外观连续变化，说明文本 token 能控制视觉分布。需要检查失败样本、prompt paraphrase、object count 与 spatial relation，避免只看最成功的 grid。概念组合也可能来自训练数据共现，而非抽象规则。
\end{knowledgebox}

\subsection{多样文本到图像样本}

接下来比较多组 prompt 与 samples，观察 diversity 与 instruction adherence 是否同时成立，并记录 gallery selection 可能带来的偏差。同一行回答“模型能生成多少种合理候选”，跨行则回答“改变文本条件是否可预测地改变对象、材质和风格”。

\lecturefigure{slide-26-dalle-examples.jpg}{多组 DALL-E samples 展示对象、材质、风格和文字条件的组合。}{官方视频 30:13。}

\begin{knowledgebox}{读图：一行内看 diversity，跨行看 instruction adherence}
同一 prompt 的多样样本显示 distribution breadth；不同 prompt 的对比显示条件是否被遵循。还要记录每个 grid 从多少候选中筛选、是否使用 reranker，以及失败率。没有 selection protocol，gallery 无法估计真实用户体验。
\end{knowledgebox}

\subsection{Zero-shot image-to-image}

进一步把条件生成扩展到 image-to-image。模型可把输入图像重描述或变换到新概念；这里的 zero-shot 指未针对该具体 transformation 训练专用 mapping，但 pretraining data 可能包含相似关系。

\lecturefigure{slide-27-dalle-zero-shot-image2image.jpg}{DALL-E zero-shot transformation 把视觉输入与文本条件组合。}{官方视频 30:33。}

\begin{knowledgebox}{读图：观察内容保持与属性改变}
理想 transformation 保留身份/布局，同时改变 prompt 指定属性。两者常冲突：改变过少是不服从，改变过多是 identity drift。评价需要 content preservation、condition adherence 与 perceptual quality 三类 metric，而非只看美观。
\end{knowledgebox}

\lecturefigure{slide-28-dalle-more-image2image.jpg}{更多 transformation examples 暴露组合泛化的范围与不稳定性。}{官方视频 31:33。}

\begin{knowledgebox}{读图：多例的价值是寻找边界}
跨多个 source/prompt 观察哪些对象、姿态和背景容易保留，哪些发生模式坍缩或属性泄漏。少数成功样本支持“能力存在”，不能估计“能力可靠”。工程上应随机抽样、盲评并报告拒绝/失败分布。
\end{knowledgebox}

\begin{warningbox}{视觉生成中的“像”与“对”}
图像没有单一逐像素正确答案，但仍有事实、身份、版权、偏见和安全约束。Perceptual quality 不能替代 provenance、consent 与 policy evaluation；生成模型的开放输出空间使评估更难，而不是无需评估。
\end{warningbox}

\subsection{本章小结}

DALL-E 把 autoregressive pretraining 扩展为 multimodal conditional generation。Gallery 显示能力存在，可靠性则需要随机采样、组合测试和明确 selection protocol。

